\documentclass[11pt]{article}
\usepackage[margin=1in]{geometry}
\usepackage{amsmath,amssymb,mathtools}
\usepackage{fontspec,unicode-math}
\setmainfont{Latin Modern Roman}
\setsansfont{Latin Modern Sans}
\setmathfont{Latin Modern Math}
\usepackage{xurl,hyperref}
\hypersetup{colorlinks=true,urlcolor=blue,linkcolor=blue}
\newcommand{\sref}[2]{\hyperlink{source:#1:#2}{[S#2]}}
\newcommand{\cataloguescope}[1]{\paragraph{Recorded mathematical question.}#1}
\begin{document}
\section{Brumer–Stark conjecture}\label{statement}

\subsection{Definitions and mathematical statement}

\paragraph{Shared notation.}$\mathbb N=\{1,2,\ldots\}$, and $\mathbb Z,\mathbb Q,\mathbb R,\mathbb C$ denote the integers, rationals, reals and complex numbers. Any other conventions are specified locally.


Let $K/k$ be a finite abelian Galois CM-extension of number fields: $k$ is totally real and $K$ is a totally imaginary quadratic extension of a totally real field. Let $G=\operatorname{Gal}(K/k)$ and let $j\in G$ be complex conjugation. Let $S$ contain all infinite places and all finite places ramified in $K/k$. Write $w_K=|\mu(K)|$ for the number of roots of unity in $K$ and
\[\theta_{K/k,S}=\sum_{\chi\in\widehat G}L_S(0,\overline\chi)\,e_\chi,\qquad e_\chi=\frac1{|G|}\sum_{g\in G}\chi(g^{-1})g.\]
Here $L_S(s,\chi)$ is the Artin $L$-function with Euler factors at $S$ removed. For an integral group-ring element $a=\sum_g a_g g$ and fractional ideal $\mathfrak A$, put $\mathfrak A^a=\prod_g g(\mathfrak A)^{a_g}$. The standard integrality theorem gives $w_K\theta_{K/k,S}\in\mathbb Z[G]$. An anti-unit $\alpha\in K^\times$ means $\alpha\,j(\alpha)=1$.
\paragraph{Statement recorded in the supplied source.}
For every fractional ideal $\mathfrak A$ of $K$, prove that an anti-unit $\alpha\in K^\times$ exists with
\[\mathfrak A^{w_K\theta_{K/k,S}}=(\alpha),\qquad K(\alpha^{1/w_K})/k\text{ abelian}.\]
This is the integral Brumer--Stark assertion, including its $2$-primary part. In particular, the predicted group-ring element annihilates the ideal class of $\mathfrak A$. \sref{statement}{2}
\subsection{Short English statement}
Prove the full Brumer--Stark ideal principalization and abelian root-extension assertion, including the prime two.
\subsection{Sources}
\begin{description}
\item[\hypertarget{source:statement:1}{S1}] ULAM.AI and the UnsolvedMath Contributors, \emph{UnsolvedMath: A Curated Collection of Open Mathematics Problems}, catalogue entry. CC BY 4.0. \url{https://huggingface.co/datasets/ulamai/UnsolvedMath}.
\item[\hypertarget{source:statement:2}{S2}] Jiro Nomura, \emph{On non-abelian Brumer and Brumer-Stark conjectures for monomial CM-extensions}, arXiv:1307.1279v1 (2013), Section 1, Conjecture 1.2, pp. 1--2; Section 2.3, Stickelberger definition, p. 9. \url{https://arxiv.org/pdf/1307.1279}.
\end{description}
\label{end:statement}
\end{document}
